
\section{计算机}

\subsection{冯诺依曼电子计算机}

冯诺依曼结构的五个主要组成部分：
\begin{enumerate}
  \item 运算器
  
  首先计算机要有运算处理数据的能力，所以需要一个处理单元来完成各种算数运算和逻辑运算，这就是算术逻辑单元（Arithmetic Logic Unit，ALU）。ALU的主要功能就是在控制信号的作用下，完成加、减、乘、除等算术运算以及与、或、非、异或等逻辑运算以及移位、补位等运算。

  运算器的主要部件就是ALU，运算器的处理对象是数据，所以数据的长度以及数据的表示方法，对运算器的影响很大。大多数通用计算机是以16、32、64位数据作为运算器一次处理数据的长度。能够对一个数据的所有位同时处理运算器称为并行运算器，一次只能对数据的一个位处理的运算器称为串行运算器。

  计算机运算时，运算器的操作对象和操作种类由控制器决定。运算器操作的数据从内存中读取，处理的结果再写入内存（或者暂时存放在内部寄存器中），而且运算器对内存数据的读写是由控制器来进行的。
  \item 控制器
  
  控制器又称为控制单元（Control Unit），是计算机的神经中枢和指挥中心，只有在控制器的控制下，整个计算机才能够有条不紊地工作、自动执行程序。

    控制器的工作流程为：从内存中取指令、翻译指令、分析指令，然后根据指令的内存向有关部件发送控制命令，控制相关部件执行指令所包含的操作。

    控制器和运算器共同组成中央处理器（Central Processing Unit），CPU是一块超大规模集成电路，是计算机运算核心和控制核心，CPU的主要功能是解释计算机指令以及处理数据。
  \item 存储器
  
  存储器的主要功能是存储程序和各种数据，并且能够在计算机运行过程高速、自动地完成程序或者数据的存储，存储器是有记忆的设备，而且采用俩种稳定状态的物理器件来记录存储信息，所以计算机中的程序和数据都要转换为二进制代码才可以存储和操作。

  存储器可以分为内部存储器（内存）和外部存储器，俩者在计算机系统中各有用处，下面大概介绍一下俩种存储器的特点：

  \begin{itemize}
    \item 内部存储器
    
    内部存储器称为内存或者主存，是用来存放欲执行的程序和数据。

    在计算机内部，程序和数据都是以二进制代码的形式存储的，它们均以字节为单位（8位）存储在存储器中，一个字节占用一个存储单元，并且每个存储单元都有唯一的地址号。

    CPU可以直接使用指令对内部存储器按照地址进行读写俩种操作，读：将内存中某个存储单元的内容读出，送入CPU的某个寄存器中；写：在控制器的控制下，将CPU中某寄存器内容传到某个存储单元中。

    内存按工作方式的不同又可以分为俩部分：

    RAM：随机存储器，可以被CPU随机读取，一般存放CPU将要执行的程序、数据，断电丢失数据

    ROM：只读存储器，只能被CPU读，不能轻易被CPU写，用来存放永久性的程序和数据，比如：系统引导程序、监控程序等。具有掉电非易失性。
    \item 外部存储器
    
    外部存储器主要来存放”暂时“用不着的程序和数据，可以和内存交换数据。

    一般是磁盘、光盘、U盘、硬盘等。
  \end{itemize}
  \item 输入设备
  \item 输出设备
  
    鼠标键盘是输入设备、显示器是输出设备；

    手机触摸屏即时输入设备又是输出设备；

    服务器中网卡既是输入设备又是输出设备；

    所有的计算机程序都可以抽象为输入设备读取信息，通过CPU来执行存储在存储器中的程序，结果通过输出设备反馈给用户。
\end{enumerate}

\section{计算机中的信息表示}


\subsection{整型}

在现代操作系统中，int 一般占用 4 个字节（Byte）的内存，共计 32 位（Bit）。如果不考虑正负数，当所有的位都为 1 时它的值最大，为 $2^{32}-1 = 4,294,967,295 ≈ 43$亿，这是一个很大的数，实际开发中很少用到，而诸如 1、99、12098 等较小的数使用频率反而较高。

使用 4 个字节保存较小的整数绰绰有余，会空闲出两三个字节来，这些字节就白白浪费掉了，不能再被其他数据使用。现在个人电脑的内存都比较大了，配置低的也有 2G，浪费一些内存不会带来明显的损失；而在C语言被发明的早期，或者在单片机和嵌入式系统中，内存都是非常稀缺的资源，所有的程序都在尽力节省内存。

反过来说，43 亿虽然已经很大，但要表示全球人口数量还是不够，必须要让整数占用更多的内存，才能表示更大的值，比如占用 6 个字节或者 8 个字节。

让整数占用更少的内存可以在 int 前边加 short，让整数占用更多的内存可以在 int 前边加 long。

\subsection{整型的长度}
我们在描述 short、int、long 类型的长度时，只对 short 使用肯定的说法，而对 int、long 使用了“一般”或者“可能”等不确定的说法。这种描述的言外之意是，只有 short 的长度是确定的，是两个字节，而 int 和 long 的长度无法确定，在不同的环境下有不同的表现。

实际情况也确实如此，C语言并没有严格规定 short、int、long 的长度，只做了宽泛的限制：
\begin{itemize}
  \item short 至少占用 2 个字节。
  \item int 建议为一个机器字长。32 位环境下机器字长为 4 字节，64 位环境下机器字长为 8 字节。
  \item short 的长度不能大于 int，long 的长度不能小于 int。
\end{itemize}

目前我们使用较多的PC系统为 Win XP、Win 7、Win 8、Win 10、Mac OS、Linux，在这些系统中，short 和 int 的长度都是固定的，分别为 2 和 4，大家可以放心使用，只有 long 的长度在 Win64 和类 Unix 系统下会有所不同，使用时要注意移植性。

\begin{table}[htpb]
  \centering
  \begin{tabular}{@{}lccc@{}}
  \toprule
  {类型}       & {16位系统/字节} & {32位系统/字节} & {64位系统/字节} \\ \midrule
  {short}     & {2}        & {2}        & {2}        \\
  {int}       & {2}        & {4}        & {4}        \\
  {long}      & {4}        & {4}        & {8}        \\
  {long long} & {8}        & {8}        & {8}        \\ \bottomrule
  \end{tabular}
  \end{table}

\subsection{定点数}

在现实生活中，我们经常使用整数和小数，不知道你有没有思考过，这些数字在计算机中是如何存储的？

如果我们想在计算机中，既能表示整数，也能表示小数，关键就在于这个小数点如何表示？

于是人们想出一种方法，即约定计算机中小数点的位置，且这个位置固定不变，小数点前、后的数字，分别用二进制表示，然后组合起来就可以把这个数字在计算机中存储起来，这种表示方式叫做「定点」表示法，用这种方法表示的数字叫做「定点数」。

定点数如果要表示整数或小数，分为以下三种情况：

\begin{enumerate}
  \item 纯整数：例如整数100，小数点其实在最后一位，所以忽略不写
  \item 纯小数：例如：0.123，小数点固定在最高位
  \item 整数+小数：例如1.24、10.34，小数点在指定某个位置
\end{enumerate}

以 1 个字节（8 bit）为例，我们可以约定前 5 位表示整数部分，后 3 位表示小数部分。

对于数字 1.5 用定点数表示就是这样：

$$1.5(\mathrm{D}) = 00001 100(\mathrm{B})$$

数字 25.125 用定点数表示就是这样：

$$25.125(\mathrm{D}) = 11001 001(\mathrm{B})$$

但是有没有发现一个问题，我们约定了前 5 位表示整数部分，后 3 位表示小数部分，此时这个整数部分的二进制最大值只能是 11111，即十进制的 31，小数部分的二进制最大只能表示 0.111，即十进制的 0.875。

如果我们想要表示更大范围的值，怎么办？

\begin{itemize}
  \item 扩大 bit 的宽度：例如使用 2 个字节、4 个字节，这样整数部分和小数部分宽度增加，表示范围也就变大了
  \item 改变小数点的位置：小数点向后移动，整个数字范围就会扩大，但是小数部分的精度就会越来越低，没有办法表示类似 0.00001 这种高精度的值
\end{itemize}

由此我们发现，不管如何约定小数点的位置，都会存在以下问题:

\begin{itemize}
  \item 数值的表示范围有限（小数点越靠左，整个数值范围越小）
  \item 数值的精度范围有限（小数点越靠右，数值精度越低）
\end{itemize}

总的来说，就是用定点数表示的小数，不仅数值的范围表示有限，而且其精度也很低。要想解决这 2 个问题，所以人们就提出了使用「浮点数」的方式表示数字。

\subsection{浮点数}

浮点数是采用科学记数法的方式来表示的，例如十进制小数 8.345，用科学记数法表示，可以有多种方式：
\begin{align*}
  8.345 &= 8.345 \times 10^0 \\
  8.345 &= 83.45 \times 10^{-1} \\
  8.345 &= 834.5 \times 10^{-2}
\end{align*}
看到了吗？用这种科学记数法的方式表示小数时，小数点的位置就变得「漂浮不定」了，这就是相对于定点数，浮点数名字的由来。

使用同样的规则，对于二进制数，我们也可以用科学记数法表示，也就是说把基数 10 换成 2 即可。

浮点数的格式可以写成这样：
\[
  \mathrm{V} = (-1)^\mathrm{S} \times \mathrm{M} \times \mathrm{R}^\mathrm{E}
\]
其中各个变量的含义如下：
\begin{itemize}
  \item S：符号位，取值 0 或 1，决定一个数字的符号，0 表示正，1 表示负
  \item M：尾数，用小数表示
  \item R：基数，表示十进制数 R 就是 10，表示二进制数 R 就是 2
  \item E：指数，用整数表示
\end{itemize}

如果我们要在计算机中，用浮点数表示一个数字，只需要确认这几个变量即可。

\subsection{浮点数标准}

1985年，IEEE 组织推出了浮点数标准，就是我们经常听到的 IEEE754 浮点数标准，这个标准统一了浮点数的表示形式，并提供了 2 种浮点格式：

\begin{itemize}
  \item 单精度浮点数 float：32 位，符号位 S 占 1 bit，指数 E 占 8 bit，尾数 M 占 23 bit
  \item 双精度浮点数 double float：64 位，符号位 S 占 1 bit，指数 E 占 11 bit，尾数 M 占 52 bit
\end{itemize}

为了使其表示的数字范围、精度最大化，浮点数标准还对指数和尾数进行了规定：

\begin{itemize}
  \item 尾数 M 的第一位总是 1（因为 $1 \leq M < 2$），因此这个 1 可以省略不写，它是个隐藏位，这样单精度 23 位尾数可以表示了 24 位有效数字，双精度 52 位尾数可以表示 53 位有效数字
  \item 指数 E 是个无符号整数，表示 float 时，一共占 8 bit，所以它的取值范围为 $0 \sim 255$。但因为指数可以是负的，所以规定在存入 E 时在它原本的值加上一个中间数 127，这样 E 的取值范围为 $-127 \sim 128$。表示 double 时，一共占 11 bit，存入 E 时加上中间数 1023，这样取值范围为 $-1023 \sim 1024$。
\end{itemize}

除了规定尾数和指数位，还做了以下规定：

\begin{itemize}
  \item 指数 E 非全 0 且非全 1：规格化数字，按上面的规则正常计算
  \item 指数 E 全 0：非规格化数，尾数隐藏位不再是 1，而是 0(M = 0.xxxxx)，这样可以表示 0 和很小的数
  \item 指数 E 全 1，尾数全 0：正无穷大/负无穷大（正负取决于 S 符号位）
  \item 指数 E 全 1，尾数非 0：NaN(Not a Number)。一些运算的结果不为实数或无穷时，就会返回这样的 NaN 值，比如计算 {\color{Green} $\sqrt{-1}$} 或 {\color{Green} $\infty -\infty$ } 时。
\end{itemize}

\subsection{阶码为什么用无符号整数表示？}

如果阶码（指数）也用补码来表示，就会使得一个浮点数中出现两个符号位：浮点数自身的和浮点数指数部分的。这样的结果是，在比较两个浮点数大小时，无法像比较整数时一样使用简单的无逻辑的二进制比较。

故而浮点数的指数部分采用了移码（无符号整数）来表示。

\subsection{浮点数为什么有精度损失？}

如果我们现在想用浮点数表示 0.2，它的结果会是多少呢？

0.2 转换为二进制数的过程为，不断乘以 2，直到不存在小数为止，在这个计算过程中，得到的整数部分从上到下排列就是二进制的结果。

\begin{lstlisting}
  0.2 * 2 = 0.4 -> 0
  0.4 * 2 = 0.8 -> 0
  0.8 * 2 = 1.6 -> 1
  0.6 * 2 = 1.2 -> 1
  0.2 * 2 = 0.4 -> 0（发生循环）
  ...
\end{lstlisting}

所以 0.2(D) = 0.00110...(B)。

因为十进制的 0.2 无法精确转换成二进制小数，而计算机在表示一个数字时，宽度是有限的，无限循环的小数存储在计算机时，只能被截断，所以就会导致小数精度发生损失的情况。

\subsection{浮点数的范围和精度有多大？}

以单精度浮点数 float 为例，它能表示的最大二进制数为 $+1.111111...1 \times 2 ^{127}$（小数点后23个1），而二进制 $1.11111...1 ≈ 2$，所以 float 能表示的最大数为 $2^{128} = 3.4 \times 10^{38}$，即 float 的表示范围为：$-3.4 \times 10^{38} \sim 3.4 \times 10^{38}$。

它能表示的精度有多小呢？

float 能表示的最小二进制数为 0.0000....1（小数点后22个0，1个1），用十进制数表示就是 $(\frac{1}{2})^{23}$。

用同样的方法可以算出，double 能表示的最大二进制数为 +1.111...111（小数点后52个1） $\times 2^{1023} ≈ 2^{1024} = 1.79 \times 10^{308}$，所以 double 能表示范围为：$-1.79 \times 10^{308} \sim 1.79 \times 10^{308}$。

double 的最小精度为：0.0000...1(51个0，1个1)，用十进制表示就是 $(\frac{1}{2})^{52}$。

从这里可以看出，虽然浮点数的范围和精度也有限，但其范围和精度都已非常之大，所以在计算机中，对于小数的表示我们通常会使用浮点数来存储。
\subsection{浮点数小结}
\begin{enumerate}
  \item 浮点数一般用科学记数法表示
  \item 把科学记数法中的变量，填充到固定 bit 中，即是浮点数的结果
  \item 在浮点数提出的早期，各个计算机厂商各自制定自己的浮点数规则，导致不同厂商对于同一个数字的浮点数表示各不相同，在计算时还需要先进行转换才能进行计算
  \item 后来 IEEE 组织提出了浮点数的标准，统一了浮点数的格式，并规定了单精度浮点数 float 和双精度浮点数 double，从此以后各个计算机厂商统一了浮点数的格式，一直延续至今
  \item 浮点数在表示小数时，由于十进制小数在转换为二进制时，存在无法精确转换的情况，而在固定 bit 的计算机中存储时会被截断，所以浮点数表示小数可能存在精度损失
  \item 浮点数在表示一个数字时，其范围和精度非常大，所以我们平时使用的小数，在计算机中通常用浮点数来存储
\end{enumerate}

\section{字符编码}


\subsection{GB2312}

GB2312兼容ASCII，采用两个值>=128的字节表示一个汉字，也就是说，它是一个变长编码，解析原则是：

\begin{itemize}
  \item 如果一个字节在0~127，那么表示一个普通ASCII字符
  \item 如果一个字节在128~255，那么继续读下一个字符（且下一个字符是128~255），两个字符联合起来的值代表一个汉字
\end{itemize}

\subsection{GBK}

GBK于 1995 年发布，将字符数扩充至21919个字符。

GBK在编码上沿袭了GB2312的做法，兼容ASCII，用双字节表示中文，但是从字符数量我们很容易看出，它不可能遵循“双字节都是128~255”这一规范，因为这个范围只有128，双字节能表示的则是128×128=16384个字符，不可能在20000+的数量

所以，既要兼容ASCII，又要限定非ASCII字符是双字节，那只能在第二个字节做妥协了，也就是说，GBK中的相当多的非ASCII字符，是第一个字符128~255，而第二个字符是0~127的，而我们知道很多文本处理是把文件当字节序列来看待，这就可能出现一些问题。

“珅”这个字在GBK编码中，第二个字符是竖线，干扰了日志格式

当然，双字节的第二个字节是0~127的，一般也都是比较生僻的汉字，但它们很多时候会出现在姓名中，在文本处理上会出现麻烦

要避免类似的问题，统一用utf8就好了，为什么utf8可以规避类似问题呢，看看utf8的编码规定就知道了。

\subsection{UTF8}

作为一种多字节变长编码，utf8通过每个字节的前缀bit严格区分了它们的作用：

\begin{itemize}
  \item 高位是0的，为单字节的ASCII字符
  \item 高位10的，是多字节序列中的后续字节
  \item 高位是11开头的，是多字节序列的首字节
\end{itemize}

这些规定不但使得不同字符的编码字节不会互相混淆，也确定了多字节序列的边界，抗乱码能力也非常强，丢失单字节是不会造成乱码的，很容易纠错，只有丢失了连续多个字节的时候，才可能出现错误内容。

\subsection{区位码、国标码、汉字机内码}

在网络传输的过程中：A要给B传输一个数据，如果A传去的数据是用GB2312-80编的，而B读取地方式是用ASCII方式的，GB 2312-80将区位码传过去，当B收到时会先读到区码，因为区码和位码的范围用十进制表示都是$0\sim 93$，如果此时的区码的十进制是6，那么ASCII对应的6则是ACK，会发现ACK是回包的指示。那么就会出现问题。

为了解决这一问题，科学家就将区位码加上32，因为ASCII的通信字符是$0\sim 32$，如果将区位码的区码和位码都+32，那么区码和位码的范围就变成了$32\sim 125$。就避开了与ASCII中通信字符与控制字符的相同的情况，提高了ASCII和区位码的兼容性。而这种$32\sim 125$的区位码就是所谓的：国标码（GB2312 -80）。

国标码是用十六进制表示的，实际上是再区位码的区码和位码分别加上20H(32O)，也就是加上2020H，最后得到了GB 2312-80。

国标码是为了防止与ASCII的控制字符和通信字符冲突，但是还是会和打印字符发送冲突，因此为了彻底让国标码和ASCII兼容。就再GB 2312-80的基础上分别再加上8080H,也就是在区码和位码上再加上128(80H)。最后得到了：汉字机内码，GB 2312就是汉字机内码。

最后在计算机内部，使用的就是汉字机内码（GB 2312）

汉字机内码也有别的称呼：汉字内码、内码。

\subsection{汉字在计算机内部的输入和输出}

如果此时要用输入法打出一个“内”字，那么我需要在键盘上点4个键:nei1。

输入nei1后，输入法会将nei1转换成国标码，然后计算机系统或者是引用软件，会将国标码再转换成汉字机内码，最后计算机或是手机再去保存这个汉字机内码。       

输入法也可能会直接将nei1转换成机内码。

输出：

汉字要输出，就需要用到另一种编码：汉字字形码。

汉字字形码是通过用一块一块的像素拼组而成的，1表示这里的像素需要图上颜色，0表示不用。从而达到输出一个汉字。 

计算机会将汉字机内码转换成国标码，再将国标码转换成汉字字形码，或是直接由汉字机内码转换成汉字字形码。最后输出。 

汉字机内码转换成汉字字形码需要经过以下步骤：

\begin{enumerate}
  \item 根据汉字机内码的编码规则，将每个汉字转换成相应的二进制编码。
  \item 对于一些早期的编码规则，需要使用专门的编码转换表将汉字机内码转换成Unicode编码。
  \item 根据汉字字形库中的字形信息，将汉字机内码转换成相应的汉字字形码。
  \item 在显示设备上根据汉字字形码渲染出汉字的图形形态，供用户观看和使用。
\end{enumerate}

\section{声音、图像、视频编码}

\subsection{图像位深}

\begin{itemize}
  \item 8位：$2^8=2^3(\mathrm{R})\times 2^3(\mathrm{G})\times 2^2(\mathrm{B})=256$色
  \item 16位：$2^16=2^5(\mathrm{R})\times 2^6(\mathrm{G})\times 2^6(\mathrm{B})=65536$色
  \item 24位：$2^24=2^8(\mathrm{R})\times 2^8(\mathrm{G})\times 2^8(\mathrm{B})=16777216$色
  \item 32位：Alpha透明度 + 24位
\end{itemize}

8/16位深度：单原色(R/G/B)最大只能表示为($0 \sim 2^3$) / ($0 \sim 2^6$)，显示的颜色范围有限。

24位：每个像素点由8bit 红 8bit 绿 8bit 蓝组合颜色而成，每一个原始颜色(R/G/B)都可以完全显示(0$\sim$0xff)，所以24位及以上，我们就叫做真彩色。

RGB色彩系统可根据Red，Green和Blue颜色的组合来构造所有颜色。

红色，绿色和蓝色分别使用8位，其整数值介于0到255之间。这使256 * 256 * 256 = 16777216种可能的颜色成为可能。

LED监视器中的每个像素通过红色，绿色和蓝色LED（发光二极管）的组合以这种方式显示颜色。

当红色像素设置为0时，LED熄灭。当红色像素设置为255时，LED完全打开。

它们之间的任何值会将LED设置为部分发光。

